% Propietario conservado: /Users/ruben/Documents/New project/output/EXTREMA_MEDIA_RAZON_RECIPROCIDAD_ES_EN_20260918/ARTICULO/ES/source/nuclear/09.tex
\chapter{Retorno observable y persistencia del registro}
\label{x_reciprocidad:nuclear:9}
\section{Recuperación conjunta desde los retornos completos}

Sea \(X\) el dominio de historias completas calibradas de
\cref{x_reciprocidad:ret:extension}, incluido el origen del calendario y, cuando se
utilice, el índice discreto de la carta. Los estados conservan las
transiciones y el registro de memoria. Escribimos
\(r=\operatorname{Ret}_+:X\to Y_+\) para el muestreo en los retornos
completos, con \(Y_+=r(X)\).

El \cref{x_reciprocidad:ret:homeomorfismo} proporciona su inversa explícita:
\(x_{3j+s}=F^s(y_j)\), para \(s=0,1,2\). La reconstrucción utiliza
el estado completo conservado en cada retorno y la actualización
efectiva; \eqref{x_reciprocidad:ret:conjugacion} conserva también el transporte.

Sean \(a:X\to R_{\rm HMT}\) la evaluación numérica por cilindros y
\(q:X\to E_\infty\) la incidencia con su calibre, según
\cref{x_reciprocidad:sec:generacion,x_reciprocidad:exc:doble-lectura,x_reciprocidad:exc:limite-inverso}.
Sea \(d:X\to X\) un transporte reversible y continuo, con el calibre
transportado. Sobre \(Y_+\) se definen



\[
a_+=a\circ r^{-1},\qquad q_+=q\circ r^{-1},\qquad d_+=r\circ d\circ r^{-1}.
\]



Entonces



\begin{equation}
\boxed{(a,q)=(a_+,q_+)\circ r,\qquad rd=d_+r.}
\label{x_reciprocidad:nuc09:lecturas}
\end{equation}



\begin{proof}
Es el \cref{x_reciprocidad:ret:lecturas}: se cancelan las composiciones con \(r\)
y \(r^{-1}\), y la inversa de \(d_+\) es \(rd^{-1}r^{-1}\).
Ambos lectores se transportan desde el mismo estado completo.
\end{proof}

En particular, para \(N\in\ZZ_{>0}\),
\(d^Nx=x\) si y sólo si \(d_+^Nr(x)=r(x)\).
La conjugación conserva la periodicidad o aperiodicidad del estado
completo. La dinámica de un lector parcial se determina mediante
su propio factor de transporte.

\section{Representación posterior explícita del transporte}

Sobre el conjunto de historias ya generado se considera la realización
posterior \(H=\ell^2(X)\), con base ortonormal
\((\delta_x)_{x\in X}\) y medida de conteo. Cada vector posee soporte
a lo sumo numerable. El transporte reversible determina



\[
C\delta_x=\delta_{d x},\qquad C^{-1}\delta_x=\delta_{d^{-1}x}.
\]



La biyección \(d\) prueba que \(C\) es unitario.
Esta representación conserva las etiquetas completas de las historias.
El espacio \(X\) y su transporte preceden a esta realización lineal.

La carta nonádica de \cref{x_reciprocidad:nuclear:2} determina



\[
T=\frac{8I+C}{9},\qquad
\eta v=\frac1{27}(8(C-I)v,-(C-I)v,\ldots,-(C-I)v),
\]





\[
I-T^*T=\eta^*\eta=\frac8{81}(I-C)^*(I-C).
\]



Sea \(P\) la proyección ortogonal sobre \(\ker(I-C)\).
La prueba espectral de \cref{x_reciprocidad:nuclear:2,x_reciprocidad:nuclear:3} da
\(T^N\to P\) fuertemente. La telescopía proporciona la isometría



\[
\mathcal Vv=(Pv,\eta v,\eta Tv,\eta T^2v,\ldots),\qquad
\mathcal M=\operatorname{Im}\mathcal V.
\]



El desplazamiento de registros \(B\) conserva el primer componente
y elimina el primer bloque de memoria. Sobre la imagen compatible,



\begin{equation}
\boxed{\widehat C=\mathcal VC\mathcal V^{-1}
=(9B-8I)|_{\mathcal M}.}
\label{x_reciprocidad:nuc09:transporte-recuperado}
\end{equation}



La inversa de la representación está dada por



\begin{equation}
v=Pv+\sum_{j\ge0}T^{*j}\eta^*(\eta T^jv).
\label{x_reciprocidad:nuc09:inversa-registro}
\end{equation}



\begin{proof}
La identidad de defecto telescopa hasta \(N\); su límite conserva
exactamente la norma. Las identidades
\(\mathcal VT=B\mathcal V\) y \(C=9T-8I\) dan
\eqref{x_reciprocidad:nuc09:transporte-recuperado}.
La isometría \(\mathcal V\) es sobreyectiva sobre \(\mathcal M\);
tomar su adjunto da \eqref{x_reciprocidad:nuc09:inversa-registro}, con convergencia
en norma de la serie. La unitariedad del transporte recuperado
se establece en \(\mathcal M\).
\end{proof}

Para \(A\in\mathcal B(H)\), sea
\(\widehat A=\mathcal VA\mathcal V^{-1}\) en \(\mathcal M\).
Esta correspondencia preserva norma, productos, adjuntos y conmutadores,
pues \(\mathcal V^{-1}=\mathcal V^*|_{\mathcal M}\).
Representa conjuntamente los observables numéricos, los de incidencia
y sus relaciones con \(C\).

\section{Convergencia numérica y correlaciones con incidencia}

Los cilindros de \eqref{x_reciprocidad:eq:cilindro} proporcionan, en la carta de
índice entero \(m(x)\), seis trits por bloque. Escribimos



\[
\xi_n(x)=\sum_{k=1}^{6n}d_k(x)3^{-k},\qquad
\xi(x)=a(x)-m(x),\qquad 0\le\xi(x)-\xi_n(x)\le3^{-6n}.
\]



Aquí \(d_k\in\{0,1,2\}\) es la lectura ternaria posterior del trit;
su orientación balanceada conserva su tipado propio. La carta
retiene \(m\) por separado, incluida la convención de extremos,
y reconstruye \(a=m+\xi\). En extremos de doble representación,
\(\xi\) se interpreta en esa carta.

Sea \(M_f\) la multiplicación por \(f\) en \(\ell^2(X)\).
La normalización \(0\le\xi_n\le\xi\le1\) hace que
\(M_\xi\) y \(M_{\xi_n}\) sean contracciones.
Para un conjunto \(B\) del codominio incidencial finito de \(q_j\),
definimos
\(E_{j,B}=M_{\mathbf1_{q_j^{-1}(B)}}\).
Es el proyector sobre las historias que satisfacen esa incidencia.

\begin{theorem}[Conservación simultánea de precisión y transporte]
\label{x_reciprocidad:nuc09:precision}
Para todo \(t\in\ZZ\),



\begin{equation}
\boxed{
\left\|\widehat C^{-t}\widehat M_\xi\widehat C^t
-\widehat C^{-t}\widehat M_{\xi_n}\widehat C^t\right\|
=\|M_\xi-M_{\xi_n}\|\le729^{-n}.}
\label{x_reciprocidad:nuc09:error-observable}
\end{equation}



Las relaciones entre los proyectores incidenciales, sus productos
y su transporte se conservan exactamente.
Para un producto ordenado de contracciones que contenga \(r\)
observables numéricos \(M_\xi\), evaluados a tiempos enteros
arbitrarios, y cualquier número finito de factores incidenciales
o unitarios exactos, reemplazar esos \(r\) factores por
\(M_{\xi_n}\) cambia el producto en norma a lo sumo



\begin{equation}
\boxed{r\,729^{-n}.}
\label{x_reciprocidad:nuc09:error-producto}
\end{equation}



La misma cota vale para cualquier elemento de matriz entre vectores
unitarios. Los factores pueden ser no conmutativos; se conserva su orden.
\end{theorem}

\begin{proof}
La norma de un multiplicador sobre \(\ell^2(X)\) es
\(\sup_{x\in X}|f(x)|\); la cola ternaria prueba la primera cota.
Conjugar por \(C^t\) y por \(\mathcal V\) conserva la norma.
Para productos \(A_1\cdots A_s\) y \(B_1\cdots B_s\), se usa



\[
\prod_{i=1}^s A_i-\prod_{i=1}^s B_i
=\sum_{j=1}^s B_1\cdots B_{j-1}(A_j-B_j)A_{j+1}\cdots A_s.
\]



Los factores exactos dan diferencia cero; cada uno de los \(r\)
restantes aporta norma a lo sumo \(729^{-n}\).
La submultiplicatividad y la cota uno de los demás factores prueban
\eqref{x_reciprocidad:nuc09:error-producto}.
Cauchy--Schwarz da la afirmación sobre elementos de matriz.
\end{proof}

Esta uniformidad conserva la precisión al transportar el observable
aproximado completo. El operador \(C\) y la incidencia permanecen
exactos en \eqref{x_reciprocidad:nuc09:error-observable}.
El índice \(n\) mide profundidad de cilindro; \(t\) mide iteraciones
del transporte \(d\); \(j\), en \eqref{x_reciprocidad:nuc09:inversa-registro},
mide iteraciones de compresión.

\section{Covariancia del operador de memoria bajo retorno}

Se restringen \(X\) y su representación atómica al sector reversible
graduado \(X_{\rm rev,grad}\) del reloj de
\eqref{x_reciprocidad:eq:weyl-relojes}. Se conserva en cada historia el grado entero
de memoria \(q_{\rm mem}\). Para \(d=\Gamma_9\), su actualización es
\(q_{\rm mem}(dx)=q_{\rm mem}(x)+1\).
El operador de memoria \(D\) es la multiplicación por
\(q_{\rm mem}\), con dominio



\[
\operatorname{Dom}D=\{v:\sum_{x\in X}|q_{\rm mem}(x)|^2|v_x|^2<\infty\}.
\]



El cambio de índice \(x\mapsto dx\) y las desigualdades
\((q\pm1)^2\le2q^2+2\) prueban que \(C\) y \(C^{-1}\)
conservan ese dominio. En él,



\begin{equation}
C^*DC=D+I,\qquad
\widehat C^*\widehat D\widehat C=\widehat D+I_{\mathcal M},
\qquad\operatorname{Dom}\widehat D=\mathcal V\operatorname{Dom}D.
\label{x_reciprocidad:nuc09:memoria}
\end{equation}



\begin{proof}
Sobre \(\delta_x\), el primer miembro vale
\(q_{\rm mem}(dx)\delta_x=(q_{\rm mem}(x)+1)\delta_x\).
La identidad se extiende al dominio indicado por truncación del
soporte en la norma del grafo de \(D\).
La segunda igualdad se obtiene por conjugación unitaria hacia
\(\mathcal M\).
\end{proof}

Por inducción,
\(q_{\rm mem}(d^Nx)=q_{\rm mem}(x)+N\).
Cada entero \(N>0\) cambia el estado completo.
Toda órbita de \(d\) es infinita; una función cuadrado-sumable
constante sobre cada una de esas órbitas es cero.
En esta realización bilateral \(P=0\).
La fórmula \eqref{x_reciprocidad:nuc09:inversa-registro} reconstruye entonces
el estado lineal completo desde sus complementos de memoria.
Esta conclusión utiliza la graduación efectiva del reloj.

El retorno de fase coexiste con el incremento de memoria de
\eqref{x_reciprocidad:nuc09:memoria}.
La aperiodicidad esturmiana y su complejidad pertenecen al factor
simbólico de \cref{x_reciprocidad:sec:generacion}; el presente resultado conserva
la dinámica graduada del estado completo.

\section{Incidencia, refinamiento y tres regímenes}

El registro \(K\) posee además la realización completa de
\cref{x_reciprocidad:nuclear:5},
\(\mathcal Fk=(\mu(k),IP_0k/6)\).
Su inversa conserva las doce coordenadas, incluida la media.
Aplicar \(\mathcal F\) a las fibras correspondientes y conservar
las etiquetas de historia transporta esta representación.
La aplicación \(\mathcal F\circ K\) conserva exactamente las fibras
del lector \(K\); la recuperación de historias completas procede
de \(r\) o de sus etiquetas efectivas.

Los sistemas de refinamiento isométrico de
\cref{x_reciprocidad:nuclear:1,x_reciprocidad:nuclear:3,x_reciprocidad:nuclear:5} conservan
\(\mathcal V\) y \(C\) mediante sus cuadrados de entrelazamiento.
La prueba utiliza las realizaciones por nivel allí construidas.
La representación atómica \(\ell^2(X)\) y las realizaciones
\(L^2\) de medida proyectiva conservan sus correspondientes
dominios y medidas.

El enlace con pliegue y despliegue se desarrolla en
\cref{x_reciprocidad:euler-trit-generadores,x_reciprocidad:euler-trit-evaluaciones}:
Coxeter, transporte nilpotente y acción de \(F_{\rm av}^2\).
Para cada realización bidimensional \(E\) que conserva la forma
alternante \(\Omega\), la carta nonádica satisface



\begin{equation}
\Omega=T_E^*\Omega T_E+D_E^*\Omega D_E,
\qquad T_E=(8I+E)/9,\quad D_E=\sqrt8(E-I)/9.
\label{x_reciprocidad:nuc09:balance-alternante}
\end{equation}



La expansión algebraica cancela los términos cruzados y utiliza
\(E^*\Omega E=\Omega\). En la realización real,
\(*\) significa transpuesta.
La ley orientada reúne los tres regímenes mediante las formas
y transportes de \cref{x_reciprocidad:nuclear:7,x_reciprocidad:nuclear:8}.
Su realización como acción utiliza además las normalizaciones
y conexiones allí especificadas.

La relación con la dinámica de conexión se formula mediante
las construcciones de \cref{x_reciprocidad:nuclear:6,x_reciprocidad:nuclear:7,x_reciprocidad:nuclear:8}.
Las identidades
\eqref{x_reciprocidad:nuc09:transporte-recuperado}--\eqref{x_reciprocidad:nuc09:balance-alternante}
conservan su representación, la precisión de los observables
y las formas transportadas.
